% Generated draft: requires mathematical and source review.
\documentclass[11pt]{article}
\usepackage[margin=25mm]{geometry}
\usepackage{fontspec}
\setmainfont{Latin Modern Roman}
\usepackage{amsmath,amssymb,mathtools,unicode-math}
\setmathfont{Latin Modern Math}
\usepackage{enumitem,needspace,xurl,hyperref}
\setlength{\parindent}{0pt}
\setlength{\parskip}{5pt}
\setlength{\emergencystretch}{3em}
\newcommand{\N}{\mathbb{N}}
\newcommand{\Z}{\mathbb{Z}}
\newcommand{\Q}{\mathbb{Q}}
\newcommand{\R}{\mathbb{R}}
\newcommand{\C}{\mathbb{C}}
\newcommand{\F}{\mathbb{F}}
\newcommand{\A}{\mathbb{A}}
\newcommand{\PP}{\mathbb{P}}
\newcommand{\E}{\mathbb{E}}
\newcommand{\eps}{\varepsilon}
\newcommand{\dd}{\,\mathrm{d}}
\begin{document}
\section*{UnsolvedMath ID: 800662}
\textbf{Source problem code:} catalogue entry
\subsection{Definitions and mathematical statement}
The supplied record does not specify the class of spaces under consideration or its conventions for admissible covers. The following gives a conditional formulation using conventional cover definitions; these conventions require confirmation against the source.

Let \(\mathscr K\) be the intended class of topological spaces. For each \(X\in\mathscr K\), let \(\mathcal C(X)\) be the intended class of admissible open covers of \(X\), with \(X\notin\mathcal U\) for every \(\mathcal U\in\mathcal C(X)\). Any countability restriction on these covers remains unspecified. Put \(\N=\{1,2,\ldots\}\). For \(n\in\N\), define
\[
\mathcal O_n(X)=\left\{\mathcal U\in\mathcal C(X):\text{every }F\subseteq X\text{ with }|F|\le n\text{ is contained in some }U\in\mathcal U\right\},
\]
and define the family of \(\omega\)-covers by
\[
\Omega(X)=\left\{\mathcal U\in\mathcal C(X):\text{every finite }F\subseteq X\text{ is contained in some }U\in\mathcal U\right\}.
\]
A cover \(\mathcal U\) is large if every point of \(X\) belongs to infinitely many distinct members of \(\mathcal U\). Define \(\mathrm T(X)\) to consist of the large covers \(\mathcal U\in\mathcal C(X)\) such that, for all \(x,y\in X\), at least one of
\[
\{U\in\mathcal U:x\in U,\ y\notin U\},\qquad
\{U\in\mathcal U:y\in U,\ x\notin U\}
\]
is finite. These are the \(\tau\)-covers under the stated convention.

For families \(\mathcal A_n\) of covers and a family \(\mathcal B\) of covers, the property \(\mathsf S_1(\{\mathcal A_n\}_{n\in\N},\mathcal B)\) means that for every sequence \((\mathcal U_n)_{n\in\N}\) with \(\mathcal U_n\in\mathcal A_n\), there exist choices \(U_n\in\mathcal U_n\) such that \(\{U_n:n\in\N\}\in\mathcal B\). Write \(\mathsf S_1(\mathcal A,\mathcal B)\) when \(\mathcal A_n=\mathcal A\) for every \(n\).

Set
\[
P(X):=\mathsf S_1(\{\mathcal O_n(X)\}_{n\in\N},\mathrm T(X)),\qquad
Q(X):=\mathsf S_1(\Omega(X),\mathrm T(X)).
\]
The supplied Conjecture 35 asserts that \(P\) is strictly stronger than \(Q\), meaning
\[
\bigl(\forall X\in\mathscr K,\ P(X)\Longrightarrow Q(X)\bigr)
\quad\text{and}\quad
\bigl(\exists X\in\mathscr K,\ Q(X)\land\neg P(X)\bigr).
\]
The intended class \(\mathscr K\) and admissible-cover conventions must be recovered before this conditional formulation can be treated as an exact source statement.
\subsection{Short English statement}
The conjecture says that choosing one member from each successive \(n\)-cover to obtain a \(\tau\)-cover is strictly stronger than making such choices from a sequence of \(\omega\)-covers. Strictness requires a space satisfying the second property but failing the first. The intended spaces and cover conventions are unspecified in the supplied record.
\subsection{Sources}
\begin{itemize}
\item Local UnsolvedMath registry: permanent ID \(800662\), exact source problem code \texttt{catalogue entry}. The supplied record labels the problem open; this is a supplied status, not an independently verified mathematical or literature assessment.
\item Tsaban, \emph{Selection Principles in Mathematics: A Milestone of Open Problems} (2003), Conjecture 35. Supplied source URL: \url{https://arxiv.org/abs/math/0312182}. The record lists Tsaban as proposer and 2003 as the proposed year.
\item The source-specific attribution/context is \texttt{[strongdiags, Conjecture 1]}. The supplied record does not resolve this bibliography label into a full reference or identify its authors.
\item Supplied provenance: access date 2026-10-16; extraction method \texttt{arxiv\_latex}. These are supplied provenance statements; no independent source consultation is claimed.
\end{itemize}
\end{document}
